% Chapter nc_v3_intro: Introduction of Volume III, the path question. Sources: Nick's draft misc/reviews/nick-volume3/vol3/ch1_intro.tex (3-23, 52-76, 78-95, 97-126, 128-154, 177-225, 227-302, 304-315, 317-326) and preface.tex, rewritten for the merged order; referee U01; 17 Nov 2026.
\label{v3:ch:intro}
The two earlier volumes built a machine. Volume I gave the theory: the L\'evy and Thorin representations of the Riemann xi function, the $2\sin^2$ lemma that carries the one to the other, and the reciprocal xi function as an extended generalised gamma convolution. Volume II turned to computation: the Thorin measure at level one, its descent into the critical strip by Cauchy deconvolution, and the numerical methods that go with it. This volume asks one question of that machine, and tries to answer it as carefully as we can.

\section{The question}\label{v3:sec:question}

Is there a path to the Riemann Hypothesis from the Thorin programme?

The question has a precise form. The Thorin measure at level $\alpha\ge1$ is positive, and it is computed from the primes. At the centre of the strip it is $\sum_{\gamma>0}\delta_{\gamma^2}$ if and only if RH holds. A path is a proof that positivity survives the descent from level one to level $\frac12$, or a proof of one of the statements equivalent to that positivity. The book already contains several such statements. We list the ones that matter here, with the labels at which they are proved.

The target in its density form is Proposition~\ref{pol:prop:pf}. Let $Y$ be the reciprocal Xi law and $Z$ a standard normal variable independent of $Y$. Then $X=Z\sqrt{2Y}$ has a smooth symmetric density $f_X$ with characteristic function $\xi(\frac12)/\xi(\frac12+t)$, and RH holds if and only if $f_X$ is a P\'olya frequency function of every order, that is, if and only if every determinant $\det[f_X(x_i-y_j)]$ with $x_1<\dots<x_k$ and $y_1<\dots<y_k$ is nonnegative. The conditions for $k=1,2,3,\dots$ form a ladder. Rung one says $f_X\ge0$. Rung two says $f_X$ is log concave. Throughout this volume we write $F$ for $f_X$ and call it the Riemann density.

The target in its book forms is Section~\ref{route:sec:target}, where six statements are proved equivalent to RH. Among them are the descent in level (the minimum of the Thorin density at the origin persists from level one to level $\frac12$), Cauchy deconvolution to the line through positive measures, Weil positivity, and the ferromagnetic form of P\'olya's law. We do not restate them here. We say only that each is a positivity statement about one explicit object built from the gamma factor and the primes, and that each is equivalent to RH, so that a proof of any one of them is a proof of RH.

A statement of the form ``the path exists'' would be a proof. A statement of the form ``the path does not exist'' cannot be proved, because RH may be true. What can be done is to ask which roads are closed, how far each open road has been walked, and what exactly is missing at its end. This is what the volume does.

\section{A guide to the six parts}\label{v3:sec:guide}

The chapters are arranged in six parts. Each chapter opens with a short note that says what it contains and what it takes from the earlier volumes. The paragraphs below say what each chapter does and why it stands where it stands.

\subsection*{Part I. The question and the tools}

\paragraph{Routes, filters and the target} This chapter combines material moved from Volume I with new sections. It assembles the results of that volume into one route to RH, states which steps are theorems, and states the one that is missing (Section~\ref{route:sec:target} opens it). It compares the route with the one setting in which RH is a theorem, curves over finite fields, where the missing step is supplied by geometry. It also states three filters that rule out whole classes of argument. It then sets out the path in detail (from Section~\ref{d:sec:kinds} on): what kind of statement is missing, how much room there is, the requirements any proof along it must meet, the three routes weighed one by one, and the milestones for which either outcome is informative. It ends with a work plan.

\paragraph{The toolkit and the cosine (Chapter~\ref{ch:cosine})} The objects and the short proofs of the facts on which everything rests, and then a solved model. The cosine takes the place of $\Xi$. Every quantity has a closed form, and a quadruple of zeros planted off the line shows what each instrument does when RH fails. The chapter also introduces the Davenport and Heilbronn control (Section~\ref{v3:sec:habits}).

\subsection*{Part II. Levels: the Thorin measure across the strip}

\paragraph{Nakamura's measure (Chapter~\ref{n:ch:nakamura})} Nakamura showed that $\xi(\sigma-it)/\xi(\sigma)$ is a characteristic function whose L\'evy measure is signed. The chapter shows that one measure serves every level, and that its only negative mass is the prime number theorem density. The criteria of P\'olya and Gneiting, which would give positivity from convexity, are defeated by the oscillation of the prime term.

\paragraph{The gamma factor and the primes (Chapter~\ref{g:ch:gammaprimes})} For that measure to be the measure of an entire function with the right zeros, several cancellations must occur, some carried by the gamma factor, one on average by the primes, and a remainder made of the zeros. The chapter asks whether the gamma function has an Euler product behind this, and brings in Schoenberg and Curry: the law of the Riemann density is a limit of splines whose knots are the reciprocal zeros. Higher splines give a family of lemmas of the $2\sin^2$ type.

\paragraph{Bondesson's letter (Chapter~\ref{bo:ch:bondesson})} A letter of Bondesson proves that $e^{-\gamma\sqrt u}/\Gamma(1+\sqrt u)$ is a generalised gamma convolution transform. The chapter takes the letter as far as it goes: the gamma block in closed form, the level criterion, the archimedean and arithmetic parts of the Thorin measure, and the descent from the levels $\frac34$ and $\frac78$. Where Volumes I and II already contain a statement, the chapter keeps the text, credits the source and cites it by label.

\paragraph{The ledger and the oscillator (Chapter~\ref{bo:ch:ledger})} The cancellation ledger of the explicit formula, the Fej\'er criterion that makes every zero pay positively (Section~\ref{bo:sec:fejerledger}), and the oscillator, a function of the cutoff alone whose values lie in a band of explicit width under RH. It also treats the value distribution of the ledger off the critical line (Section~\ref{bo:sec:values}).

\paragraph{Families and level one (Chapter~\ref{bo:ch:families})} What the oscillator can do beyond $\zeta$ and beyond the critical line: families of $L$ functions, the level one inequality with its proof from the primes, the control, the step from the half plane $\frac78$ of \cite{OAI26a} to a level just below it, and a Dirichlet character of conductor $12$.

\paragraph{The classical sources (Chapter~\ref{ch:classics})} The oscillator has ancestors. The chapter reads four of them: von Mangoldt's proof of Riemann's formula, Ramanujan on highly composite numbers, the long paper of Hardy and Littlewood with the criterion of Riesz, and the book of Hardy, Littlewood and P\'olya. Each is read in the language of the volume, and what it contributes is separated from what it does not.

\paragraph{Chebyshev's error and the ladder (Chapter~\ref{ct:ch:cheb})} Chebyshev's two memoirs and the statistics of $\psi(x)-x$ read as Thorin measures. The 1850 identity is a ratio of Gamma functions, and it is shown to be a generalised gamma convolution (Theorem~\ref{ct:thm:chebggc}). The chapter also gives the periodogram of the normalised error and a sufficient condition for a signed product of Gamma functions to be a generalised gamma convolution.

\paragraph{Iwaniec's approach} This chapter is moved from Volume II (Section~\ref{iw:sec:approach} opens it). It reports the approach of Iwaniec to the level of the Thorin measure, and the dictionary between his objects and ours. The chapters on Bondesson cite it for the rank and origin statements.

\paragraph{The quasi Riemann hypothesis} One moved chapter, in three parts, reads the zero free half planes for Dirichlet and Hecke $L$ functions against the instruments of the volume. Section~\ref{qrh:sec:claims} states what is proved, by whom, and with what status. The half plane $\re s>\frac78$ is the theorem of \cite{OAI26a}.

\subsection*{Part III. Total positivity and the heat flow}

\paragraph{The P\'olya frequency ladder (Chapter~\ref{p:ch:path})} The central chapter of the part. For $x>\pi/4$ the Riemann density is a residue expansion whose signs alternate exactly when consecutive zeros are simple and on the line, and rung $k$ holds in this far field whenever the first $k$ zeros are simple and on the line. The difficulty sits in the near field. The chapter proves that $F$ is strictly log concave (Theorem~\ref{p:thm:m1}), by ball arithmetic. It then proves the detection rung theorem: if the zeros below the lowest zero off the line are $m$ simple zeros on the line and that zero is one simple quadruple, the ladder fails at rung $m+1$. A finite set of rungs therefore sees only finitely many zeros, and the content of RH lies in a mechanism that proves all rungs at once. It embeds $F$ in the de Bruijn and Newman heat flow.

\paragraph{Kernels, splines and Hankel ladders (Chapter~\ref{c:ch:cand})} Two candidate mechanisms for that step are tested against the control. Euler factors are totally positive on their own lattices, and what they build is total positivity of Riemann's kernel, which stops at a finite order and does not see the zeros. Spline knots built from primes converge to the ordinates of the zeros whether or not the zeros lie on the line. Both fail in their naive form. The chapter then identifies the Wald couple that carries $\xi$ to $1/\xi$, gives two more ladders computable from the Taylor coefficients of $\log\xi$, and arrives at the Thorin string of $\xi$.

\subsection*{Part IV. Chains}

\paragraph{The birth and death chain (Chapter~\ref{b:ch:chain})} The Stieltjes continued fraction of the Thorin measure at the centre is computed from the moments of P\'olya's kernel, and its coefficients are read as the rates of a birth and death chain. RH holds if and only if every rate is positive. The rates are conjectured to follow a closed law of Bohr and Sommerfeld type (Conjecture~\ref{b:conj:bs}), carried by the gamma factor, multiplied by an arithmetic fluctuation of a few per cent. RH is the statement that the fluctuation never reaches $-1$. For the control the chain fails at exactly the levels its zeros off the line predict.

\paragraph{The Bohr and Sommerfeld law (Chapter~\ref{l:ch:bohr})} The law for the gamma chain has two halves. The chapter proves the direct half: the chain defined by the law has the counting function of the smooth zeros. The inverse half is not proved, and the chapter says what is needed.

\paragraph{Chains at every level (Chapter~\ref{h:ch:chains})} The chain at the centre is placed at the bottom of a family indexed by a level $\alpha>\frac12$. The rates at level $\alpha$ are all positive if no zero has real part $\beta\ge\alpha$, and, under the proviso of Theorem~\ref{h:thm:crit}, only if no zero has real part $\beta>\alpha$. Above $\alpha=1$ the rung is the Euler product, at $\alpha=1$ it is the theorem of Hadamard and de la Vall\'ee Poussin, and RH is the rung at $\frac12$. In every computed case each zero to the right of the level costs two negative rates (Conjecture~\ref{h:conj:count}).

\paragraph{Linear response of the rates (Chapter~\ref{ch:progress})} An exact formula for the response of every rate to every zero, which shows which zeros carry the arithmetic fluctuation; an exact determinant formula for a pair of zeros off the line; and a collision experiment in which two zeros approach each other and then leave the line. This is the first step along the route of the rates.

\paragraph{The Berry Keating operator} This chapter is moved from Volume I (Section~\ref{ch:bk} opens it). It is placed after the chains because the Krein string of the zeros is what the rates would be rates of.

\subsection*{Part V. Spins}

\paragraph{Spins and instruments (Chapter~\ref{f:ch:ferro})} P\'olya's density read as a spin in the sense of Lee and Yang. A ferromagnetic law is a hitting time in disguise, so P\'olya's approach and the Thorin programme are one question. The chapter also collects what Ramanujan's identities and the Vere-Jones laws contribute as instruments, each with its status.

\subsection*{Part VI. What is forced, and is there a path}

\paragraph{What the Euler product and the functional equation force} This chapter is moved from Volume II (Section~\ref{rg:sec:target} opens it). It asks what the two structural properties of $\zeta$ force on the zeros, and ends in one sign condition that sits next to the Fej\'er criterion of Chapter~\ref{bo:ch:ledger}.

\paragraph{Is there a path? (Chapter~\ref{chap:conclusions})} The closing chapter gathers the equivalent forms of the target, the roads that are closed, the record of the control on every test, and the open problems of this volume in one numbered list (Section~\ref{fin:sec:open}).

\section{What is proved in this volume}\label{v3:sec:proved}

The status words are those of the book: proved, computed (floating point), certified (ball arithmetic only), conditional on a named hypothesis, heuristic, open. Each item below carries one of them. The chapter that proves it is named.
\begin{itemize}
\item The Riemann density $F$ is strictly log concave on the real line, hence strictly positive (Theorem~\ref{p:thm:m1}, Chapter~\ref{p:ch:path}). Status: certified. This settles the PF$_2$ part of Problem~\ref{pol:op:pf2} of Volume I, the second rung of the ladder, with every decision a comparison of balls.
\item If the zeros below the lowest zero off the line are $m$ simple zeros on the line and that zero is one simple quadruple, rung $m+1$ of the ladder fails (Theorem~\ref{p:thm:detect}, Chapter~\ref{p:ch:path}). Status: proved. For the Davenport and Heilbronn control, with $43$ zeros on the line below the first quadruple, the theorem predicts, under hypotheses checked only numerically, that fixed configurations of up to $43$ points pass in the far field and that rung $44$ fails (computed).
\item The smooth window theorem: under RH the last cutoff at which the smoothed ledger at level $\frac12+\delta$ is negative, multiplied by $\sqrt\delta$, tends to an explicit constant as $\delta\to0$ (Theorem~\ref{bo:thm:smooth}, Chapter~\ref{bo:ch:ledger}). Status: conditional on RH, with an effective bracket.
\item The Fej\'er ledger holds at every height. For every length $L>0$ the Ces\`aro truncated Thorin density is a sum over zeros of kernels that are nonnegative whenever the level lies to the right of the zero, minus an explicit gamma term (Theorem~\ref{bo:thm:cesaro}, Section~\ref{bo:sec:fejerledger}). Status: proved, with no condition on the cutoff.
\item The value distribution of the ledger at a level $\alpha\in(\frac12,1)$ has a limit law with a smooth density and support the whole line (Theorem~\ref{bo:thm:limitlaw}, Section~\ref{bo:sec:values}). Status: proved, with no hypothesis. The limit law is due to Kershner and Wintner~\cite{KershnerWintner37}, and the chapter says what it adds.
\item Chebyshev's Gamma ratio is a generalised gamma convolution (Theorem~\ref{ct:thm:chebggc}, Chapter~\ref{ct:ch:cheb}). Status: proved.
\item The direct half of the Bohr and Sommerfeld law (Chapter~\ref{l:ch:bohr}). Status: proved, to leading order, with the error term of Theorem~\ref{l:thm:law}. Its numerical agreement is computed.
\item Positivity of the Thorin measure of the reciprocal Xi function at every level $\alpha>\frac78$, and the GGC property of the compensated function on each such rung (Theorem~\ref{dz:thm:qrh}, Section~\ref{dz:sec:qrh}). Status: proved, from the theorem of \cite{OAI26a}.
\item Under RH the oscillator takes its values in a band of width $0.09242$ (Theorem~\ref{bo:thm:osc}, Chapter~\ref{bo:ch:ledger}). Status: conditional on RH.
\end{itemize}
What is conditional, and on what. Statements marked RH assume the Riemann Hypothesis and are used only to read what RH would imply. Statements marked GRH assume it for the $L$ functions named. Statements marked LI assume the linear independence of the ordinates of the zeros over the rationals. The half plane $\re s>\frac78$ for Dirichlet and Hecke $L$ functions is proved in \cite{OAI26a}, and Section~\ref{qrh:sec:claims} records its status. Statements below $\frac78$ that need a zero free half plane stay conditional on a named hypothesis $(\mathrm Q_\theta)$.

\section{What is not proved}\label{v3:sec:notproved}

Nothing in this volume proves the Riemann Hypothesis. The step below level $\frac78$ is open: no argument in the volume shows that the positivity of the Thorin measure at level $\frac78$, which follows from \cite{OAI26a}, survives to any level in $(\frac12,\frac78)$ without an input of the strength of a zero free half plane. What the volume offers is a target stated exactly, in forms that can be computed and tested, a set of theorems and certified computations around it, and a list of closed roads. The finite tests, however long, see only finitely many zeros, as the detection rung theorem shows. A proof would have to supply a mechanism that holds at every order at once.

\section{Two habits}\label{v3:sec:habits}

Two habits run through the volume.

The first is a control. The Davenport and Heilbronn function has a functional equation of the same shape as $\xi$, a Dirichlet series and a zero counting function of the same shape, no Euler product, and zeros off the critical line. Every test we build is run on it as well. A test that the control passes is not a test of RH, and an argument that would apply equally to the control is wrong. The record of the control is kept in the closing chapter.

The second is a ledger of evidence. Each statement is marked with one of the words listed in Section~\ref{v3:sec:proved}, and one convention serves the whole volume. Floating point computations are marked computed, with the precision where it matters. Only ball arithmetic is called certified. A heuristic is a plausibility argument and nothing more. We have tried never to let one status pass for another.

\section{How the volume was written}\label{v3:sec:written}

The volume grew from notes of Nicholas Polson, written in November 2026 as a sequence of chapters, each beginning from a question about the one before. The notes were refereed chapter by chapter, and their numbers were recomputed. Where a number could not be recomputed it is marked. Statements that are already in Volumes I and II are cited by label and not repeated, except in the chapters on Bondesson, which keep the whole text with credit. Results that are new in a chapter are credited there to an unrefereed note of Polson (November 2026). We make no claim of priority beyond what a chapter says. The normalisation is the book's throughout: $u=w^2$, atoms at $\gamma^2$, and $\xi$ with Riemann's factor $\frac12$.
